\documentclass{article}

\usepackage{amsmath,amssymb}
\usepackage{xurl}
\usepackage[hidelinks]{hyperref}

% Shared commands for notes in docs/paper-gaps/.
% Notation follows the LDT paper (references/ldt-paper/).

\newcommand{\Lean}{\textsc{Lean}}
\newcommand{\leanid}[1]{\textsf{#1}}
\newcommand{\ghissue}[1]{\href{https://github.com/LionSR/MIPStarRE/issues/#1}{GitHub issue \##1}}
\newcommand{\ghpr}[1]{\href{https://github.com/LionSR/MIPStarRE/pull/#1}{GitHub PR \##1}}

% --- Paper-compatible notation ---
\newcommand{\F}{\mathbb{F}}
\newcommand{\N}{\mathbb{N}}
\newcommand{\C}{\mathbb{C}}
\newcommand{\tr}{\operatorname{tr}}
\newcommand{\ot}{\otimes}
\newcommand{\E}{\mathop{\bf E\/}}
\newcommand{\bx}{\mathbf{x}}
\newcommand{\by}{\mathbf{y}}
\newcommand{\bu}{\mathbf{u}}
\newcommand{\bv}{\mathbf{v}}

% Approximation relation (paper uses \approx_\eps and \simeq_\zeta)
\newcommand{\approxeps}{\approx_{\varepsilon}}
\newcommand{\simgeq}{\simeq_{\zeta}}

% Polynomial spaces (paper uses \polyfunc, \polysub, \polymeas)
\newcommand{\polyfunc}[3]{\operatorname{Poly}(#1,#2,#3)}
\newcommand{\polysub}[3]{\operatorname{PolySub}(#1,#2,#3)}
\newcommand{\polymeas}[3]{\operatorname{PolyMeas}(#1,#2,#3)}


\title{Scalar approximation chain for \texttt{lem:comm-data-processed-g}}
\author{}
\date{\today}

\begin{document}
\maketitle

\section{The assertion}

This note concerns the proof of \Cref{lem:comm-data-processed-g} in
\cite{Ji2020LowIndividualDegree}
(\texttt{commutativity-G.tex}, lines~72--131). The lemma asserts
that, under the hypotheses of consistency, strong self-consistency,
and boundedness for a family $\{G^x\}$ of projective sub-measurements,
the second term in the commutativity expansion (\texttt{eq:gcom8}) is close
to the first:
\begin{multline*}
  \Bigl|
    \E_{\bu,\bv,\bx,\by} \sum_{a,b}
      \bigl\langle \psi \bigm|
        G^{\bu,\bx}_a G^{\bv,\by}_b G^{\bu,\bx}_a G^{\bv,\by}_b
          \otimes I
      \bigm| \psi \bigr\rangle
  \\
  {} -
    \E_{\bu,\bv,\bx,\by} \sum_{a,b}
      \bigl\langle \psi \bigm|
        G^{\bu,\bx}_a G^{\bv,\by}_b G^{\bu,\bx}_a
          \otimes I
      \bigm| \psi \bigr\rangle
  \Bigr|
  \le 48m(\sqrt\gamma + \sqrt\zeta).
\end{multline*}

\section{The paper's proof route}

The paper assembles this bound through a ten-step scalar approximation
chain. The chain is:
\begin{enumerate}
  \item Insert Bob's measurement $A^{\bv,\by}_b$ at the right (cost $2\sqrt\zeta$).
  \item Remove trailing $G^{\by}$ via \Cref{clm:g-comm-stability} (cost $\sqrt\zeta$).
  \item Insert Bob's second measurement $A^{\bu,\bx}_a$ at the right (cost $2\sqrt\zeta$).
  \item Swap the two Bob measurements via \Cref{thm:commutativity-points}
    (cost $6\sqrt{\gamma(m+1)}$).
  \item Remove trailing $G^{\bx}$ via \Cref{clm:g-comm-stability2}
    (cost $\sqrt\zeta + 6\sqrt{\gamma(m+1)}$).
  \item--7. Reverse the two measurement insertions (cost $2\sqrt\zeta + 2\sqrt\zeta$).
  \item--9. Apply postprocessed self-consistency twice (cost $\sqrt\zeta + \sqrt\zeta$).
\end{enumerate}
The error tally is
\[
  2(12\sqrt\zeta + 12\sqrt{\gamma(m+1)})
  = 24(\sqrt{\gamma(m+1)} + \sqrt\zeta)
  \le 48m(\sqrt\gamma + \sqrt\zeta).
\]

\section{The Lean proof route}

The Lean formalization implements the same chain in
\texttt{ProcessedG.lean}\footnote{\texttt{MIPStarRE/LDT/Commutativity/ScalarApproximation/ProcessedG.lean},
  lines~1194--1300.} and organizes it into four phases:
\begin{description}
  \item[Phase~1] (steps~1--2): error $2\sqrt\zeta + \sqrt\zeta$.
  \item[Phase~2] (steps~3--5): error $2\sqrt\zeta + 6\sqrt{\gamma(m+1)} + \sqrt\zeta + 6\sqrt{\gamma(m+1)}$.
    The $6\sqrt{\gamma(m+1)}$ swap is via \texttt{commutativityPoints};
    the boundedness part of \Cref{clm:g-comm-stability2} is packaged as
    \texttt{gCommStabilityTwo\_scalar}.
  \item[Phase~3] (steps~6--7): error $2\sqrt\zeta + 2\sqrt\zeta$.
  \item[Phase~4] (steps~8--9): error $\sqrt\zeta + \sqrt\zeta$.
\end{description}
The total error is identical:
\[
  2 \cdot (12\sqrt\zeta + 12\sqrt{\gamma(m+1)})
  \le \texttt{commDataProcessedGError}
  = 48 m (\sqrt\gamma + \sqrt\zeta).
\]

The Lean proof uses the same mathematical ingredients as the paper:
\texttt{closenessOfIP} (\Cref{prop:closeness-of-ip}) for inserting and
removing measurement factors,
\texttt{commutativityPoints} (\Cref{thm:commutativity-points}) for the
point-measurement swap, \texttt{gCommStability\_scalar}
(\Cref{clm:g-comm-stability}) for the first stability claim, and
\texttt{gCommStabilityTwo\_scalar} (\Cref{clm:g-comm-stability2}) for
the second. The four-phase organization is a regrouping of the ten paper
steps, not a different proof route.

The internal splitting of \Cref{clm:g-comm-stability2} into a
\texttt{commutativityPoints} swap ($6\sqrt{\gamma(m+1)}$) and a separate
boundedness part ($\sqrt\zeta$) matches the paper's proof of that
claim (lines~192--222), where the first step uses
\Cref{thm:commutativity-points} to swap the order of $A^{\bu,\bx}_a$ and
$A^{\bv,\by}_b$.

\section{Conclusion}

The Lean scalar chain is mathematically equivalent to the
paper's ten-step chain. The same approximations are applied, the same
error constants appear, and the final bound $48m(\sqrt\gamma + \sqrt\zeta)$
is identical. The four-phase reorganization is an internal proof-structuring
choice that does not alter any mathematical assertion.

\noindent\textbf{Verdict:} no discrepancy. The Lean proof follows the
paper's scalar chain faithfully, with the same error budget and the same
intermediate theorems.

\bibliographystyle{alpha}
\bibliography{docs/paper-gaps/references}

\end{document}